\section{Positive reordering in a quantum torus}
\label{alg:section}

We prove the algebraic statement that supplies positivity at a wall joint.
The proof keeps two gradings distinct: polynomial degree records powers of
$v$, whereas an auxiliary filtration makes a coproduct regular.  The
resulting algebra need not be commutative.  Its graded vector-space
dimensions nevertheless have the required symmetric-algebra form.

Let $\Gamma$ be a lattice with an integral alternating form $\Omega$, and
let $C\subset\Gamma$ be an additive monoid contained in a strictly convex
cone in a real two-dimensional subspace.  Fix an additive charge
$c:C\to\N$ that is positive away from zero, and assume that
$\{p\in C:c(p)\le B\}$ is finite for every $B$.  The completed quantum
torus has multiplication
\begin{equation}\label{alg:torus}
 X_pX_q=v^{\Omega(p,q)}X_{p+q}
\end{equation}
and coefficients in $\Q((v^{-1}))$.  Thus a coefficient has a bounded
upper $v$-degree and can have infinitely many negative degrees.
For $p\in C\setminus\{0\}$ and $k\in\Z$, define
\begin{equation}\label{alg:units}
 E_{p,k}=\prod_{j\ge0}(1+v^{k-2j}X_p),\qquad
 H_{p,k}=\prod_{j\ge0}(1-v^{k-2j}X_p)^{-1}.
\end{equation}
We call these \emph{elementary units}.  A list of their occurrences is
\emph{admissible} if, for each charge bound $B$, their parameters $k$ are
bounded above and only finitely many occurrences with $c(p)\le B$ have
$k\ge L$, for each $L\in\Z$.  Occurrences are counted with multiplicity.

\begin{theorem}[Positive reordering]\label{alg:reordering}
Suppose an admissible product of elementary units is written in ray
order, with $\Omega(p_i,p_j)>0$ whenever the ray of $p_i$ precedes the
distinct ray of $p_j$.  Its factorization in the opposite ray order is
again a product of elementary units, with nonnegative integral
multiplicities.  The output list is admissible.

Every coefficient at fixed exponent and fixed $v$-degree is a finite
sum of nonnegative integers, and each exponent has bounded upper
$v$-degree.  The assertion also holds modulo any fixed charge cutoff.
For the algorithmic assertion, assume that lattice operations, the form
$\Omega$ and the charge are effectively computable, that input occurrences
at bounded charge and above a given parameter threshold can be listed
effectively, and that for every $B$ one can compute an integer $A_B\ge1$
such that
\[
 |\Omega(p,q)|\le A_B
 \qquad\text{whenever }p,q\in C,\quad c(p),c(q)\le B.
\]
Then every prescribed coefficient and every bounded part of the output
unit list can be calculated by finite rational linear algebra.
\end{theorem}

We first decompose shuffle spaces by slope and use separated-center
restrictions to construct a coproduct on an auxiliary associated graded
space. Polynomiality of the restriction symbols removes negative
separation powers before evaluation at zero. A signed vector-space PBW
decomposition and primitive strings then give the positive elementary
factors. Throughout, polynomial degree with its shifts records powers of
$v$; the final energy bounds control infinite input lists.

A ray factorization with constant term one exists uniquely over
$\Q((v^{-1}))$: at charge $B$, subtract the contributions already
determined at smaller charges, and assign the remaining coefficient of
$X_p$ to the ray of $p$.  Products involving that new coefficient and a
nonconstant term have larger charge.  We shall prove that, in the order
specified in Theorem~\ref{alg:reordering}, each resulting ray factor
has the stronger form asserted there.

\subsection{Shuffle spaces and destabilizing quotients}

Begin with finitely many input occurrences, indexed by $1,\ldots,t$ in
input ray order.  Units on one ray may be ordered arbitrarily.  Set
$a_{ii}=0$ for $E_{p_i,k_i}$ and $a_{ii}=1$ for $H_{p_i,k_i}$.  For
$i\ne j$, put
\[
 a_{ji}=\max\{\Omega(p_i,p_j),0\}.
\]
For dimension vectors $d,e\in\N^t$, define
\begin{align}
 p(d)&=\sum_i d_ip_i,&
 \chi(d,e)&=\sum_i d_ie_i-\sum_{i,j}a_{ij}d_ie_j,
 \label{alg:euler}\\
 s(d)&=\tfrac12\chi(d,d),&
 u(d)&=\sum_i d_i(k_i+1-a_{ii}).
 \label{alg:shifts}
\end{align}
Then
\begin{equation}\label{alg:skew}
 \chi(d,e)-\chi(e,d)=\Omega(p(d),p(e)).
\end{equation}
Choose a linear functional $\eta$ on the plane so that
$\mu(d)=\eta(p(d))/c(p(d))$ increases in input ray order.
This is possible by intersecting the strictly convex cone with
$c=1$.  Dimension vectors of a common slope have symmetric pairings
under $\chi$.  We use $|d|_1=\sum_i d_i$ only to count variables and to
perform inductions; it is different from charge.

Let
\[
 S_d=\Q[x_{i,\alpha}:1\le i\le t,\ 1\le\alpha\le d_i]
       ^{\prod_i\mathfrak S_{d_i}},\qquad S_0=\Q,
\]
with every variable of polynomial degree one.  Variables with a fixed
$i$ form a pack.  The polynomial shuffle presentation is the cohomological Hall product
of Kontsevich--Soibelman \cite[Section~2.4, Theorem~2]{KS11}; see also
\cite[Theorem~2.2]{Ef12}.  In our normalization it is
\begin{equation}\label{alg:shuffle}
 f*g=\sum_{\text{pack shuffles}}f(x)g(y)K_{d,e}(x,y),\qquad
 K_{d,e}(x,y)=
 \prod_{i,j}\prod_{\alpha\le d_i,\,\beta\le e_j}
 (y_{j,\beta}-x_{i,\alpha})^{a_{ij}-\boldsymbol1_{i=j}}.
\end{equation}
Each shuffle chooses, in each output pack, the variables assigned to
the first input; there is no factorial normalization.  The only
possible poles are simple poles between two variables of the same
output pack.  Interchanging those two variables pairs the residues
with opposite signs.  After taking a common denominator, the
numerator is therefore divisible by every such difference.  The
differences are relatively prime irreducible polynomials, so all poles
cancel.  Consequently $f*g\in S_{d+e}$.  A simultaneous shuffle of
three input groups proves associativity: each pair of groups
contributes exactly one kernel in either parenthesization.  Polynomial
degree changes by $-\chi(d,e)$.

The destabilizing quotient below is the polynomial form of the
semistable restriction quotient in \cite[Theorem~8.1]{FR18}.  We
establish its needed properties directly for our charge-weighted slope.
For $d\ne0$, put
\begin{equation}\label{alg:quotient}
 I_d=\sum_{\substack{a+b=d,\ a,b\ne0\\\mu(a)>\mu(d)}}S_a*S_b,
 \qquad B_d=S_d/I_d,
 \qquad B_0=\Q.
\end{equation}
Every image $S_a*S_b$ is an ordinary $S_d$-submodule.  Indeed an
output symmetric polynomial, restricted to the two input packs, is a
finite sum of polynomial tensors; multiplication by it can be moved
inside the shuffle.  Thus $I_d$ is an ordinary homogeneous ideal.

The product descends to $B_d\otimes B_e\to B_{d+e}$ when
$\mu(d)=\mu(e)=\theta$.  If $a+b=d$ and $\mu(a)>\theta$, a term
$(S_a*S_b)*S_e$ has the same destabilizing first group $a$.  If
$a+b=e$ and $\mu(a)>\theta$, a term
$S_d*(S_a*S_b)=(S_d*S_a)*S_b$ has destabilizing first group $d+a$.
This proves descent in both arguments.  Give $B_d$ shifted degree
\begin{equation}\label{alg:shifted-degree}
 \deg_{\mathrm{sh}}f=\deg_{\mathrm{pol}}f+s(d).
\end{equation}
On a common slope,
$s(d+e)=s(d)+s(e)+\chi(d,e)$, so the descended product preserves the
sum of shifted degrees.

\subsection{Separated centers and ordered decomposition}

We record carefully the restriction operation used throughout the
proof.  For an ordered list $P=(r_1,\ldots,r_m)$ of nonzero dimension
vectors summing to $d$, divide each pack of $d$ into labeled groups
with multiplicities $r_\ell$, replace the variables in group $\ell$
by $x^\ell+z_\ell$, and reduce modulo $I_{r_\ell}$ in every group.
This defines
\begin{equation}\label{alg:restriction-S}
 \widetilde R_P^z:S_d\longrightarrow
       \left(\bigotimes_{\ell=1}^m B_{r_\ell}\right)
          [z_1,\ldots,z_m].
\end{equation}
The choices inside a pack do not matter because its input is
symmetric.  Common translation preserves $I_d$: all kernels in
Equation~\eqref{alg:shuffle} depend on differences.  Hence translation
also defines an automorphism of $B_d[z]$.

Here is the useful way to apply Equation~\eqref{alg:restriction-S} to
a shuffle expression.  Split each source group into subgroups assigned
to the target groups.  At distinct centers expand the kernels between
different targets in their internal variables, over
$\Q(z_1,\ldots,z_m)$.  Each coefficient is separately symmetric in
the source subgroups and is a finite sum of polynomial tensors.
The kernels within one target are left unexpanded and give its
ordinary shuffle product.  If a prefix of the source groups has
dimension $b$ inside a target of dimension $r$, with
$\mu(b)>\mu(r)$, that target contribution is in $I_r$ by
associativity and vanishes.  All these operations are legitimate in
the product of homogeneous internal-degree spaces.  Vanishing of the
Taylor expansion implies vanishing of the original polynomial
restriction, since its target ideals are homogeneous and extension of
scalars to $\Q(z_1,\ldots,z_m)$ is injective.

The slope ordering follows the Harder--Narasimhan factorization
formalism \cite[Propositions~4.8 and 4.12]{Rei03},
\cite[Sections~5.2--5.3]{KS11}.  Ordered multiplication of semistable
pieces is established in \cite[Theorem~6.2]{FR18}; the next argument
proves the required decomposition for the explicit quotients above.
Here the locator for \cite{FR18} refers to its published version.

\begin{lemma}[Ordered decomposition]\label{alg:ordered}
Choose graded vector-space sections $B_d\to S_d$ of all quotient
maps.  For every $d$, multiplication induces a vector-space
isomorphism
\begin{equation}\label{alg:ordered-sum}
 S_d\cong
 \bigoplus_{\substack{d^{(1)}+\cdots+d^{(m)}=d\\
       \mu(d^{(1)})>\cdots>\mu(d^{(m)})}}
 B_{d^{(1)}}\otimes\cdots\otimes B_{d^{(m)}},
\end{equation}
with the polynomial-degree shifts prescribed by the shuffle kernels.
\end{lemma}

\begin{proof}
Induct on $|d|_1$.  Modulo the chosen lift of $B_d$, every element is a
sum of destabilizing products.  Expand their two factors by induction.
Draw a polygon for a resulting list by joining the successive points
\[
 \left(c(p(d^{(1)}+\cdots+d^{(j)})),
       \eta(p(d^{(1)}+\cdots+d^{(j)}))\right).
\]
The original destabilizing cut lies strictly above the chord joining
the endpoints.  If adjacent slopes are nondecreasing, merge those two
factors and apply the induction hypothesis to their product.  Their
total dimension is smaller than $d$: a list of only two such factors
could not have a destabilizing cut.  The replacement polygons are
concave and lie above the chord of the merged pair, hence above the
old pair.  Equality can occur only for a straight pair, in which case
the number of segments decreases.  The polygon remains strictly above
the total chord somewhere.  There are only finitely many lists of
nonzero subdimension vectors of $d$, so this process terminates at
strictly decreasing slopes.  This proves spanning.

For independence, restrict an ordered product to targets
$r_1,\ldots,r_m$ of strictly decreasing slopes.  A surviving split has
no destabilizing source prefix in any target.  At a source cut, if
the portion in target $r_j$ has charge $b_j$, its height is at most
$b_j\mu(r_j)$.  Among all choices with fixed sum $\sum_jb_j$, this
upper bound is maximized by filling the targets in decreasing slope
order.  Thus the source polygon lies below the target polygon at
every source corner, and hence everywhere: the target polygon is
concave.  If the polygons coincide, strictness of their slopes forces
equality at each corner.  The first target must be completely filled
at the first corner, the first two at the second, and so on.
Consequently the only surviving split assigns whole source groups to
the matching whole targets.  In particular the \emph{dimension
vectors}, not only their charges and slopes, must agree.

In a finite relation choose a polygon minimal under pointwise
comparison.  Restriction to its dimension-vector list kills every
other list: a surviving source polygon would have to be smaller, or
equal with that same list.  For the matching list the restriction is
the tensor of the translated lifts, multiplied by all intergroup
kernels.  These kernels are invertible Taylor series at separated
centers.  They act on the quotient tensor because they are separately
symmetric in each whole target group.  Cancel them.  Translation is
an automorphism of each quotient, and the chosen lifts identify with
the quotient vector spaces.  Thus the original tensor in that summand
must be zero.  Applying this to a minimal nonzero summand contradicts
the proposed relation.
\end{proof}

\begin{lemma}[Restrictions on a slope]\label{alg:restrictions}
If all $r_\ell$ have the slope of $d$,
Equation~\eqref{alg:restriction-S} factors through a map
\begin{equation}\label{alg:restriction}
 R_P^z:B_d\longrightarrow
       \left(\bigotimes_\ell B_{r_\ell}\right)[z].
\end{equation}
These maps compose under refinement, with the successive centers
added.  When a same-slope product is restricted, only splits whose
nonempty subgroups have this slope survive; the restrictions of its
two factors may already be taken in their respective quotients.
\end{lemma}

\begin{proof}
Consider a destabilizing product with first dimension $a$.  If every
piece of $a$ in a target had slope at most the common slope, their sum
would too, contrary to $\mu(a)>\mu(d)$.  Some target is therefore
destabilized, and the preceding Taylor test kills the term.  This
proves descent.  Refinement is composition of polynomial substitutions;
the descended maps therefore compose as claimed.

For a product with first total dimension of the common slope, split
that dimension among equal-slope targets.  If any nonempty piece is
off the slope, some piece must lie above it, by the weighted-average
identity for slopes.  That target is destabilized.  In every remaining
split both pieces of a target have the common slope.  Internal shuffle
products descend by Equation~\eqref{alg:quotient}; multiplication by
the Taylor coefficients preserves the quotient ideals because they
are ordinary ideals.  This justifies taking all input subgroups in
their quotients before the remaining computation.
\end{proof}

\subsection{The Hilbert-series target}

The ordered decomposition gives a formula for each factor in the
reversed ray order.  We derive that formula now, so that the remaining
problem is a precise statement about graded dimensions.

The Hilbert series of the ordinary pack-symmetric polynomial space is
\begin{equation}\label{alg:ordinary-hilbert}
 \Hilb_{S_d}(t)=\prod_i\prod_{j=1}^{d_i}(1-t^j)^{-1}.
\end{equation}
Indeed monomial orbit sums are indexed by weakly decreasing exponent
lists in each pack, and their successive differences contribute
independent degrees $1,\ldots,d_i$.

The coefficient of $X_{mp}$ in $E_{p,k}$ is the generating series of
$0\le j_1<\cdots<j_m$, weighted by
$v^{km-2\sum j_r}$; subtracting $0,1,\ldots,m-1$ turns this into a
weakly increasing list.  For $H_{p,k}$ the list is weakly increasing
from the outset.  Thus those two coefficients are respectively
\[
 \frac{v^{km-m(m-1)}}{\prod_{j=1}^m(1-v^{-2j})},
 \qquad
 \frac{v^{km}}{\prod_{j=1}^m(1-v^{-2j})}.
\]
Taking the ordered torus product shows that the finite input is
\begin{equation}\label{alg:input-series}
 \sum_{d\in\N^t}X_{p(d)}
       v^{u(d)-\chi(d,d)}\Hilb_{S_d}(v^{-2}).
\end{equation}
The diagonal terms give exactly the two displayed unit coefficients;
the off-diagonal terms $-\chi(d,d)$ supply
$\sum_{i<j}d_id_j\Omega(p_i,p_j)$, which is the torus multiplication
shift.

Apply Lemma~\ref{alg:ordered}.  For two blocks $a,b$, the shuffle
degree shift is $-\chi(a,b)$.  Hence normalization by
$v^{-\chi(d,d)}$ changes their combined exponent relative to the
separate normalizations by
\[
 2\chi(a,b)-\chi(a,b)-\chi(b,a)
     =\chi(a,b)-\chi(b,a)=\Omega(p(a),p(b)).
\]
This is precisely Equation~\eqref{alg:torus}.  Equation~\eqref{alg:input-series}
therefore factors in decreasing slope order, with slope factor
\begin{equation}\label{alg:slope-factor}
 1+\sum_{\mu(d)=\theta}X_{p(d)}
          v^{u(d)-\chi(d,d)}\Hilb_{B_d}(v^{-2}),
\end{equation}
where $\Hilb_{B_d}$ uses ordinary polynomial degree.

In this formula an element of shifted degree $h$ contributes the
weight $v^{u(d)-2h}$.  It therefore suffices to show that the direct
sum of $B_0$ and the spaces $B_d$ of this slope has the graded
dimensions of a tensor product of polynomial and exterior algebras,
with the generators grouped into strings of fixed dimension $d$ and
shifted degrees $h,h+1,h+2,\ldots$.  A string of polynomial
generators contributes $H_{p(d),u(d)-2h}$; a string of exterior
generators contributes $E_{p(d),u(d)-2h}$.  We will obtain exterior
strings when $\chi(d,d)$ is odd and polynomial strings when it is
even.  These are assertions about a vector-space decomposition;
they will not require the same-slope shuffle product itself to be
commutative.

The next construction supplies this decomposition through a
coproduct.  Its auxiliary filtration makes separated restrictions
regular at coincident centers, while retaining the original shifted
degree that appears in the weight $v^{u(d)-2h}$.

\subsection{A filtration and polynomial restriction symbols}

Fix a slope $\theta$ for this subsection and the next two.  All
dimension vectors in a splitting have slope $\theta$, unless they are
zero.  For each $d\ne0$ and $w\in\tfrac12\Z$, define a decreasing
filtration by
\begin{equation}\label{alg:filtration}
 f\in\cF_wB_d
 \quad\Longleftrightarrow\quad
 \text{every term of }R_P^zf\text{ has internal degree }
       +\sum_{r\in P}s(r)\ge w\text{ for every }P.
\end{equation}
Centers have degree zero in this test; the one-block list $(d)$ is
included.  There are finitely many lists $P$ for fixed $d$.  Their
internal degrees are nonnegative, giving a finite lower bound for the
filtration.  In original polynomial degree $\ell$, the one-block test
has internal degree at most $\ell$, giving a finite upper bound
$\ell+s(d)$.  The filtration is consequently exhaustive and finite
in each original degree.  Put
$\gr_wB_d=\cF_wB_d/\cF_{w+1/2}B_d$ and give $B_0=\Q$ weight zero.

For $f\in\gr_wB_d$ and $P=(r_1,\ldots,r_m)$ define
\begin{equation}\label{alg:symbol}
 L_P(f)=
 \prod_{i<j}(z_j-z_i)^{\chi(r_i,r_j)}
          (R_P^zf)_w,
\end{equation}
where the subscript retains internal degree
$w-\sum_i s(r_i)$.  A priori this expression has rational center
coefficients.

\begin{example}[The two gradings in one variable]
\label{alg:one-variable}
Let $d_i=1$ and let all other entries of $d$ be zero.  There is no
proper splitting of $d$, so $B_d=\Q[x]$ and the only restriction is
$R_{(d)}^z f=f(x+z)$.  For $f=x^\ell$ with $\ell\ge0$, the term $z^\ell$ has
internal degree zero.  Hence every $x^\ell$ has filtration weight
$s(d)=(1-a_{ii})/2$, whereas its original shifted degree is
$\ell+s(d)$.  Its symbol is
\[
 L_{(d)}([x^\ell])=z^\ell.
\]
The auxiliary grade therefore retains every polynomial degree, even
though they all have the same filtration weight.  On these symbols,
multiplication by $x$ becomes multiplication by $z$, and differentiation
becomes differentiation in $z$.  This is the simplest instance of
the polynomial strings constructed below.
\end{example}

\begin{lemma}[Polynomial symbols and detection]\label{alg:detection}
Every $L_P(f)$ is polynomial in its centers.  The maps $L_P$ jointly
inject $\gr_wB_d$ into their polynomial targets.  For the additive
tensor filtration
\begin{equation}\label{alg:tensor-filtration}
 \cF_w(B_a\otimes B_b)
    =\sum_{u+v\ge w}\cF_uB_a\otimes\cF_vB_b,
\end{equation}
tensor products of the restriction tests characterize membership,
and tensor products $L_R\otimes L_S$ jointly detect its associated
grades.
\end{lemma}

\begin{proof}
Write $S(P)=\sum_{r\in P}s(r)$.  Merge two groups $a,b$ of $P$.
The shift sum becomes $S(P)+\chi(a,b)$.  Refine this merged group
again, giving the two groups a relative center $\delta=z_b-z_a$.
This refinement is homogeneous if $\delta$ is counted, together with
the internal variables, as degree one.  By the coarse filtration
test, every term before this refinement has degree at least
$w-S(P)-\chi(a,b)$.  Therefore a term of fine internal degree
$w-S(P)$ must have $\delta$-degree at least $-\chi(a,b)$.
When $\chi(a,b)<0$, this proves exactly the divisibility required to
cancel the denominator in Equation~\eqref{alg:symbol}.
Perform this argument for every pair.  To justify simultaneous
cancellation, choose a $\Q$-basis of the quotient tensor space and
write the polynomial in that basis.  Each scalar coefficient is
divisible by the requisite power of each diagonal difference.
Distinct differences $z_j-z_i$, for unordered pairs $\{i,j\}$, are
relatively prime irreducibles, so their product divides that
coefficient.  No domain property of a quotient algebra is needed.

Joint injectivity follows directly from
Equation~\eqref{alg:filtration}: if all leading restrictions vanish,
the representative belongs to $\cF_{w+1/2}$.  Multiplication by the
nonzero center prefactors cannot change this conclusion, as can be
seen after adjoining the independent centers as a field extension.

For completeness, a tensor satisfies all lower-bound restriction
tests exactly when it belongs to Equation~\eqref{alg:tensor-filtration}.
One implication follows by applying the defining tests to each
factor.  For the other, choose filtration complements in each
finite-dimensional original-degree space.  If a tensor violated the
bound, take its lowest nonzero total filtration grade below $w$.
Tensoring the two injective families of graded restriction maps over
$\Q$ is injective.  At a fixed total grade, contributions of
different factor grades $u,v$ cannot cancel: in tests $R,S$ their two
separate internal degrees are
$u-S(R)$ and $v-S(S)$, respectively.  Thus some tensor test detects
the violating grade.  The same argument proves joint detection on
every associated tensor grade, including after adjoining center
variables or localizing their nonzero differences.
\end{proof}

\begin{lemma}[Product symbols]\label{alg:product-symbol}
The same-slope product preserves the sum of filtration weights.
For $f\in\gr_uB_d$, $g\in\gr_vB_e$, and
$R=(r_1,\ldots,r_m)$, its leading symbols satisfy
\begin{equation}\label{alg:product-symbol-formula}
 L_R(f*g)=
 \sum_{\substack{r_i=d_i'+e_i'\\
                  \sum_i d_i'=d,\ \sum_i e_i'=e}}
 (-1)^{\sum_{i>j}\chi(d_i',e_j')}
 (*^{\otimes R})
 \bigl(L_{(d_i')}(f)\otimes L_{(e_i')}(g)\bigr).
\end{equation}
All nonzero entries in this sum have slope $\theta$; zero entries in
a list of symbols are omitted.  The operator $*^{\otimes R}$ places
the two factors belonging to each target group next to each other,
without a sign, and applies its shuffle product.
\end{lemma}

\begin{proof}
Use Lemma~\ref{alg:restrictions} and expand intergroup kernels at
separated centers.  Every nonconstant internal Taylor term increases
internal degree.  In target $i$, internal shuffling changes
polynomial degree by $-\chi(d_i',e_i')$, exactly compensated by
\[
 s(r_i)-s(d_i')-s(e_i')=\chi(d_i',e_i').
\]
This proves the filtration bound.  On its leading grade only the
scalar intergroup kernels remain.  For $i<j$, their contribution is
\[
 (z_j-z_i)^{-\chi(d_i',e_j')}
 (z_i-z_j)^{-\chi(d_j',e_i')}.
\]
Multiplication by the $L_R$ prefactor cancels these exponents against
the cross terms in $\chi(r_i,r_j)$, leaving the two input prefactors
and the sign $(-1)^{\chi(d_j',e_i')}$.  Multiplying these signs over
$i<j$ gives Equation~\eqref{alg:product-symbol-formula}.
\end{proof}

\subsection{The coproduct and its regular specialization}

Let $\mathcal H=\bigoplus_{\mu(d)=\theta}\gr B_d$, with
$\mathcal H_0=\Q$.  Its product was just constructed.  We next
construct an actual coproduct on these vector spaces.  This requires
removing negative separation powers on the associated grade before
setting the separation to zero.

\begin{lemma}[Graded coproduct]\label{alg:coproduct}
For $d=a+b$ on the fixed slope there is a map
\[
 \Delta_{a,b}:\gr B_d\longrightarrow\gr B_a\otimes\gr B_b
\]
preserving both original shifted degree and total filtration weight,
characterized by
\begin{equation}\label{alg:coproduct-symbol}
 (L_R\otimes L_S)(\Delta_{a,b}f)=L_{R,S}(f)
\end{equation}
for all lists $R,S$ of $a,b$, with independent centers and with the
lists concatenated on the right.  When a dimension is zero its
component is the usual identity tensor unit.
\end{lemma}

\begin{proof}
Take first $a,b\ne0$ and a representative
$f\in\cF_wB_d$ homogeneous of polynomial degree $\ell$.  Form the
Laurent expansion at $Z=\infty$
\begin{equation}\label{alg:separation-expansion}
 K_{a,b}(x+Z,y)^{-1}f(x+Z,y)
       =\sum_j C_j(f)Z^j
       \quad\text{in }(B_a\otimes B_b)((Z^{-1})).
\end{equation}
Each coefficient is an actual polynomial tensor before quotienting:
expand the finitely many integer powers of linear factors, and note
that a specified $Z$-power receives only finitely many binomial
choices.  The expansion is separately symmetric in the two packs.
It descends from $B_d$ because the restriction descends and the target
ideals are ordinary ideals.  Homogeneity, counting $Z$ as degree one,
gives
\begin{equation}\label{alg:coefficient-degree}
 \deg_{\mathrm{pol}}C_j(f)=\ell+\chi(a,b)-j.
\end{equation}
In particular $C_j(f)=0$ for $j>\ell+\chi(a,b)$.

Apply arbitrary refinements $R=(r_i)$ of $a$ and $S=(s_j)$ of $b$ to
each coefficient.  Their centers in the first list are shifted by
$Z$.  The internal-degree-zero part of the inverse kernel is
\begin{equation}\label{alg:inverse-scalar}
 \prod_{i,j}(z_j^S-z_i^R-Z)^{\chi(r_i,s_j)},
\end{equation}
and all its other internal terms have positive degree.  The
restriction of $f$ to the concatenated list has internal shifted
degree at least $w$, by its defining filtration tests.  Thus the
refined tests of every $C_j(f)$ have the same lower bound.  Tensor
detection in Lemma~\ref{alg:detection} proves
$C_j(f)\in\cF_w(B_a\otimes B_b)$, coefficient by coefficient.

On the weight-$w$ grade only Equation~\eqref{alg:inverse-scalar}
contributes.  Multiplying by the two separate symbol prefactors
supplies exactly the missing cross-list factors, so
\begin{equation}\label{alg:polynomial-Z}
 (L_R\otimes L_S)
       \left(\sum_j[C_j(f)]Z^j\right)
   =L_{R,S}([f])(z^R+Z,z^S).
\end{equation}
By Lemma~\ref{alg:detection}, the right side is a polynomial in all
its centers and in $Z$.  Consequently every $j<0$ has zero image
under every tensor symbol.  The same Lemma implies
$[C_j(f)]=0$ for each $j<0$.  This conclusion holds in the actual
associated graded tensor space; it is not a conclusion drawn after
a specialization of a rational kernel.

Equation~\eqref{alg:coefficient-degree} leaves only finitely many
nonnegative $Z$-powers.  We may therefore evaluate this finite tensor
polynomial at zero and define
\[
 \Delta_{a,b}[f]=[C_0(f)].
\]
Changing $f$ by $\cF_{w+1/2}$ changes every coefficient by that higher
tensor filtration, so the definition is independent of the
representative.  Setting $Z=0$ in
Equation~\eqref{alg:polynomial-Z} proves
Equation~\eqref{alg:coproduct-symbol}.  Finally the original shifted
degree of $C_0(f)$ is
\[
 \ell+\chi(a,b)+s(a)+s(b)=\ell+s(d),
\]
as required.
\end{proof}

Write $\varepsilon(a,b)=(-1)^{\chi(a,b)}$.  Symmetry of $\chi$ on
the slope makes this a symmetric bilinear interchange sign.  Give
tensor products the multiplication
\begin{equation}\label{alg:signed-tensor}
 (x\otimes y)(x'\otimes y')
       =\varepsilon(\deg y,\deg x')\,(x*x')\otimes(y*y'),
\end{equation}
where degrees in this formula are dimension vectors.

\begin{proposition}\label{alg:bialgebra}
With $\Delta=\sum_{a+b=d}\Delta_{a,b}$,
$\mathcal H$ is a connected dimension-graded bialgebra, and its
coproduct is cocommutative with the interchange sign
$\varepsilon(a,b)$.
\end{proposition}

\begin{proof}
Testing either iterated coproduct against three independent lists of
symbols gives the symbol for their concatenation.  Joint tensor
detection gives coassociativity.  Swapping two lists changes their
cross-list prefactor by
$(-1)^{\chi(a,b)}$; the underlying restriction is simply permuted.
Detection therefore gives signed cocommutativity.  The empty-list
components give the counit and the unit identities.

For multiplicativity apply
Equation~\eqref{alg:product-symbol-formula} to the two tensor factors
separately.  If the first input contributes dimensions $d_R,d_S$ to
the two lists, and the second contributes $e_R,e_S$, concatenating
the lists adds to the two separate symbol signs exactly
\[
 \sum_{i\in S,\ j\in R}\chi(d_i',e_j')=\chi(d_S,e_R).
\]
This is the interchange sign in
Equation~\eqref{alg:signed-tensor}.  Both sides therefore have the
symbols of the product restricted to the concatenation, and
detection proves multiplicativity.  Positive dimension grading and
$\mathcal H_0=\Q$ give connectedness.
\end{proof}

\subsection{Primitive vectors and polynomial strings}

The symmetric-algebra form of the graded dimensions is the
cohomological-integrality and PBW phenomenon developed in
\cite{KS11,Ef12,DMe20}.  Efimov's freeness theorem assumes a symmetric
quiver, whereas our quiver need only have symmetric Euler pairing on
the selected slope.  Davison--Meinhardt use a perverse associated
graded in \cite[Theorem~C]{DMe20}; the separated-center filtration here
has its own explicit construction.  We include the vector-space
argument needed from our bialgebra, because its conclusion does not
require commutativity of the product.
Let $\mathcal P\subset\mathcal H$ be the primitive subspace, namely
the positive-dimensional elements killed by every proper
$\Delta_{a,b}$.  By Equation~\eqref{alg:coproduct-symbol} and tensor
detection, these are exactly the elements killed by all proper
symbols $L_P$, where a proper list has at least two nonempty entries.

The convolution-logarithm proof below is the signed/braided
Eulerian-idempotent argument.  For the classical cocommutative
characteristic-zero method, see
\cite[Sections~3--4, Theorem~3.5 and Lemma~4.1]{Patras93};
the proof here supplies the bilinear-sign adaptation.

\begin{lemma}[Signed vector-space PBW]\label{alg:pbw}
Symmetrized multiplication gives an isomorphism of graded vector
spaces
\begin{equation}\label{alg:pbw-map}
 \operatorname{Sym}_{\varepsilon}(\mathcal P)
       \longrightarrow\mathcal H.
\end{equation}
Here the left side is the signed symmetric vector space: exchanging
homogeneous vectors of dimensions $a,b$ multiplies a tensor by
$\varepsilon(a,b)$.  The gradings retained are dimension, original
shifted degree, and filtration weight; tensor length is summed over.
\end{lemma}

\begin{proof}
For linear maps out of $\mathcal H$, convolution uses its coproduct
and multiplication in the target.  Let $e$ be augmentation followed
by inclusion of the unit, and put
\[
 J=\id-e,\qquad
 \mathsf P=\log^*(\id)=
      \sum_{r\ge1}\frac{(-1)^{r-1}}rJ^{*r}.
\]
On dimension $d$, a nonzero term has at most $|d|_1$ positive
dimensions, so this sum and all subsequent convolution series are
finite there.  In the target algebra $\mathcal H\otimes\mathcal H$ define
$\alpha(h)=h\otimes1$ and $\beta(h)=1\otimes h$.
Then $\alpha*\beta=\Delta$, while signed cocommutativity gives
$\beta*\alpha=\Delta$ as well.  Since $\Delta$ is multiplicative,
\[
 \Delta\circ\mathsf P
   =\log^*(\Delta)
   =\log^*(\alpha)+\log^*(\beta)
   =\mathsf P\otimes1+1\otimes\mathsf P.
\]
Thus $\mathsf P$ takes values in $\mathcal P$.  Exponentiating,
\[
 \id=e+\sum_{r\ge1}\frac1{r!}\mathsf P^{*r}.
\]
The iterated coproduct is invariant under signed permutations, and
$\mathsf P$ preserves all gradings.  This expression therefore
spans $\mathcal H$ by signed symmetrized products of primitives.

To prove independence, identify signed symmetric tensors in
characteristic zero with the invariant tensors obtained by averaging
over signed permutations.  In a finite relation among their
products, let $r$ be the largest tensor length.  Apply the iterated
coproduct with $r$ factors and project each factor to positive
dimension.  Terms of length less than $r$ vanish.  A product of $r$
primitives gives the sum of their $r!$ signed permutations; on its
symmetrized tensor this is $r!$ times that tensor.  Hence the
length-$r$ part of the relation is zero.  Descending on $r$ proves
independence.

For a homogeneous basis of $\mathcal P$, a multiset contributes one
signed symmetric tensor unless an element with self-sign $-1$
repeats, in which case it contributes zero.  This proves the stated
description of dimensions.  In particular, an odd primitive may have
a nonzero square in $\mathcal H$.  Its square is then an even
primitive, since the two middle coproduct terms cancel, and is
accounted for in the even primitive subspace.  No assertion that raw odd
squares vanish is used in Equation~\eqref{alg:pbw-map}.
\end{proof}

\begin{lemma}[Primitive strings]\label{alg:strings}
For each nonzero dimension $d$, the primitive space is a direct sum
of strings with original shifted degrees
\[
 h,h+1,h+2,\ldots,
 \qquad h-s(d)\in\N.
\]
Each degree has finite dimension.  Every vector in these strings has
self-interchange sign $(-1)^{\chi(d,d)}$.
\end{lemma}

\begin{proof}
On $S_d$, consider common differentiation and multiplication by the
variable mean,
\[
 D_d=\sum_{i,\alpha}\frac{\partial}{\partial x_{i,\alpha}},
 \qquad
 T_d f=\left(\frac1{|d|_1}\sum_{i,\alpha}x_{i,\alpha}\right)f.
\]
The translation invariance of the shuffle kernels implies
$D_d(I_d)\subset I_d$.  The ideal property implies
$T_d(I_d)\subset I_d$.  Thus both operators descend to $B_d$.
For $P=(r_1,\ldots,r_m)$, restriction gives
\begin{align*}
 R_P^z(D_df)&=\left(\sum_\ell\partial_{z_\ell}\right)R_P^zf,\\
 R_P^z(T_df)&=
 \left(\frac{\sum_\ell|r_\ell|_1z_\ell}{|d|_1}
       +\frac{\sum_{\ell,i,\alpha}x^\ell_{i,\alpha}}{|d|_1}
 \right)R_P^zf.
\end{align*}
Both preserve the filtration lower bound.  On the associated grade,
the internal-variable term in the second line raises internal degree
and disappears.  The prefactor of $L_P$ depends only on center
differences, so its common center derivative vanishes.  Consequently
\begin{align*}
 L_P(D_df)&=\left(\sum_\ell\partial_{z_\ell}\right)L_P(f),\\
 L_P(T_df)&=
       \frac{\sum_\ell|r_\ell|_1z_\ell}{|d|_1}L_P(f).
\end{align*}
They preserve the common kernel of all proper symbols, hence the
primitive space.  They satisfy $D_dT_d-T_dD_d=\id$ there; $D_d$ is
locally nilpotent, and $T_d$ raises original shifted degree by one.

Here is an explicit string decomposition.  On a space with these
operators define the finite-on-each-vector operator
\[
 \pi(f)=\sum_{j\ge0}\frac{(-1)^j}{j!}T_d^jD_d^jf.
\]
The commutator identity gives $D_d\pi=0$.  Binomial cancellation
gives
\[
 f=\sum_{j\ge0}\frac{T_d^j}{j!}\pi(D_d^jf).
\]
Choose a homogeneous basis of $\ker D_d$ in the primitive space.
Its $T_d$-iterates therefore span.  They are independent: in a finite
relation $\sum_{j=0}^rT_d^jb_j=0$ with $D_db_j=0$, applying
$D_d^r$ gives $r!b_r=0$, and induction removes the other terms.

Every component of $\mathcal H_d$ is a subquotient of $B_d$, with
original shifted degree equal to nonnegative polynomial degree plus
$s(d)$.  Thus every string begins at $h\ge s(d)$ with the indicated
integral difference.  The polynomial spaces in any fixed degree are
finite-dimensional, and the filtration is finite there, proving
the finiteness assertion.  Finally the self-sign depends only on the
dimension vector, as stated.
\end{proof}

Combining Lemmas~\ref{alg:pbw} and~\ref{alg:strings}, the shifted
Hilbert series on a slope is the Hilbert series of independent
polynomial strings when $\chi(d,d)$ is even and exterior strings
when it is odd.  Passing to the auxiliary associated grade has not
changed any original shifted graded dimension.

\subsection{Elementary factors and the infinite-list bounds}

We now finish the finite-input case of
Theorem~\ref{alg:reordering} using the slope factor
\eqref{alg:slope-factor}.  Lemmas~\ref{alg:pbw}
and~\ref{alg:strings} give exactly the polynomial and exterior
strings required above.  In shifted degree $h$, their weight is
$v^{u(d)-2h}$.

A primitive
string beginning at $h$ contributes one elementary unit with
\begin{equation}\label{alg:output-parameter}
 p=p(d),\qquad k=u(d)-2h,
\end{equation}
of type $E$ if $\chi(d,d)$ is odd and type $H$ if it is even.
These are integer parameters because $2h$ is an integer.  The
primitive-basis dimensions give nonnegative integral multiplicities.
Finiteness of each original graded piece proves that, for a fixed
finite list and bounded dimension, only finitely many output
parameters lie above any threshold.  This proves the finite-input
case.

Now allow an admissible infinite input list.  Give each occurrence
its own index and use finite-support dimension vectors $d$.  All
polynomial spaces, ideals, restrictions, filtrations, and primitive
spaces for a fixed $d$ use only the finitely many indices in its
support.  Adding other indices changes none of them.  The preceding
constructions are consequently compatible as the finite input list
increases.

We give uniform bounds to justify this passage.  Fix an integer
charge cutoff $B\ge1$.  Every active $p_i$ belongs to the finite set
$C_B=\{p\in C:0<c(p)\le B\}$, and
$D=|d|_1\le B$ whenever $c(p(d))\le B$.  Choose $K\ge0$ bounding
all their input parameters above, and choose $A\ge1$ bounding
$1$ and $|\Omega(p,q)|$ for $p,q\in C_B$.  Then
\[
 \chi(d,d)\ge-AD^2,\qquad
 u(d)\le\sum_i d_ik_i+D.
\]
Since $h\ge s(d)$, Equation~\eqref{alg:output-parameter} gives
\begin{equation}\label{alg:parameter-bound}
 k\le u(d)-\chi(d,d)
       \le\sum_i d_ik_i+B+AB^2
       \le BK+B+AB^2.
\end{equation}
If an output parameter satisfies $k\ge L$, every active occurrence
has
\begin{equation}\label{alg:active-input-bound}
 k_i\ge L-(B-1)K-B-AB^2.
\end{equation}
To see this, single out one of its $d_i$ copies in the preceding
sum and bound the other at most $B-1$ copies by $K$.  Admissibility
leaves only finitely many possible input occurrences in
Equation~\eqref{alg:active-input-bound}, hence only finitely many
possible $d$.  For each one, $k\ge L$ also bounds
$h\le(u(d)-L)/2$, so only finitely many primitive starts can occur.
This proves both parts of admissibility for the output list.

Finally, consider a coefficient at exponent of charge at most $B$
and fixed energy $v^e$ in any of the unit products under discussion.
Expand its selected units into the index lists used above.  There
are at most $B$ selected monomial copies.  Let $K'\ge0$ bound the
parameters of that list at this charge, and let $A'$ bound absolute
pairings of its possible exponents.  A term has energy
\begin{equation}\label{alg:energy-bound}
 e=\sum_{\nu=1}^{D}k_\nu-2\sum_{\nu=1}^{D}j_\nu+Q,
 \qquad D\le B,\quad j_\nu\ge0,\quad |Q|\le A'B^2.
\end{equation}
It follows that
\[
 e\le BK'+A'B^2,\qquad
 k_\nu\ge e-(B-1)K'-A'B^2,
 \qquad
 \sum_\nu j_\nu\le\frac{BK'+A'B^2-e}{2}.
\]
Only finitely many unit occurrences meet the lower parameter bound,
and only finitely many nonnegative index lists meet the last bound.
Thus the coefficient is a finite nonnegative integer and its upper
energy is bounded.  These bounds justify all coefficientwise limits
of the finite-list identities, proving the infinite case.

For effectiveness, compute $K$ as the maximum of zero and the parameters
in the finite input list at charge at most $B$ and threshold zero, and
take $A=A_B$ from the supplied pairing bound.  A charge and
output-parameter cutoff now give the finite set of relevant input occurrences by
Equation~\eqref{alg:active-input-bound}.  For each resulting $d$,
the bound on $h$ gives finitely many polynomial degrees.  In those
degrees the shuffle images, quotient ideals, restriction maps,
filtrations, coproduct, primitive kernels, and $D_d$-kernels are
matrices over $\Q$ on finite polynomial spaces.  Their ranks give
the required unit multiplicities.  Equation~\eqref{alg:energy-bound}
then computes any desired product coefficient from a finite list.

The effective bound is available in the wall applications below.
There the roots lie in an explicit finite free cone with basis
$r_1,\ldots,r_s$ and coordinate-sum charge.  If $c(p),c(q)\le B$, then
\[
 |\Omega(p,q)|\le B^2\max_{i,j}|\Omega(r_i,r_j)|.
\]
Taking the maximum of this integer and one supplies $A_B$, also for
the monoids in the two-dimensional joint planes.  Thus those applications
satisfy the additional effective hypothesis.  This completes the proof
of Theorem~\ref{alg:reordering}.
